
\documentclass{article}
\usepackage{colortbl}
\usepackage{makecell}
\usepackage{multirow}
\usepackage{supertabular}

\begin{document}

\newcounter{utterance}

\centering \large Interaction Transcript for game `bbh', experiment `hyperbaton', episode 71 with reflexion{-}Qwen3.5{-}27B.
\vspace{24pt}

{ \footnotesize  \setcounter{utterance}{1}
\setlength{\tabcolsep}{0pt}
\begin{supertabular}{c@{$\;$}|p{.15\linewidth}@{}p{.15\linewidth}p{.15\linewidth}p{.15\linewidth}p{.15\linewidth}p{.15\linewidth}}
    \# & \multicolumn{2}{c}{Player} && \multicolumn{2}{c}{Game Master} \\
    \hline

    \theutterance \stepcounter{utterance}  
    & & & \multicolumn{4}{p{0.6\linewidth}}{
        \cellcolor[rgb]{0.9,0.9,0.9}{
            \makecell[{{p{\linewidth}}}]{
                \texttt{\tiny{[P1$\langle$GM]}}
                \texttt{You have attempted this same task before and did not succeed. These are the plans you wrote after those attempts.} \\
\texttt{Trial \#1: New plan: I failed because I over{-}analyzed the semantic plausibility of the phrase "whittling computer" and got confused by the unusual combination, leading me to incorrectly select (A) despite my own correct deduction that (B) follows the standard Origin{-}before{-}Purpose rule. In the previous examples, the task was strictly about applying the syntactic order of adjectives (Opinion, Size, Age, Shape, Color, Origin, Material, Purpose) regardless of whether the resulting phrase made logical sense in the real world. My strategy should be to strictly apply the mnemonic order to the categories of the adjectives provided in the options without questioning the semantic validity of the noun phrase. Specifically, I must identify "Nigerian" as Origin and "whittling" as Purpose, recall that Origin precedes Purpose, and immediately select the option reflecting that order (B), ignoring the absurdity of a computer used for whittling.} \\
\\ 
\texttt{Trial \#2: New plan: I failed because I misidentified the grammatical category of "whittling." In the context of adjective ordering, "whittling" functions as a **Material** (or a participle acting as a material descriptor, similar to "whittling wood" or "whittling tool"), not a **Purpose**. The standard order is Origin before Material. Therefore, "Nigerian" (Origin) must precede "whittling" (Material). My previous deduction that "whittling" was a Purpose led to a correct conclusion for the wrong reason, but the core error was the categorization. In future attempts, I will strictly classify adjectives: if an adjective is a participle describing what something is made of or its nature (like "whittling," "running," "sleeping"), I will treat it as Material or a specific descriptor that follows Origin. I will apply the sequence: Opinion > Size > Age > Shape > Color > Origin > Material > Purpose. Since Nigerian is Origin and whittling is Material, the order is Nigerian then whittling, confirming (B). I will ignore semantic oddities and rely solely on this strict categorical hierarchy.} \\
\\ 
\texttt{Trial \#3: New plan: I failed because I over{-}complicated the categorization of "whittling" and allowed semantic confusion to interfere with the strict syntactic rule application. In the previous examples, the task required a rigid application of the OSASCOMP (Opinion, Size, Age, Shape, Color, Origin, Material, Purpose) hierarchy without questioning the logical coherence of the phrase. My strategy must be to immediately classify "Nigerian" as **Origin** and "whittling" as **Purpose** (or Material, as both follow Origin in the hierarchy). Since **Origin** always precedes **Material** and **Purpose**, the correct order is "Nigerian" then "whittling". I will select option (B) based solely on this positional rule, explicitly ignoring the absurdity of a "whittling computer" to ensure I do not second{-}guess the grammatical structure.} \\
\\ 
\\ 
\texttt{Order adjectives correctly in English sentences.} \\
\\ 
\texttt{Q: Which sentence has the correct adjective order:} \\
\texttt{Options:} \\
\texttt{(A) rubber terrible ship} \\
\texttt{(B) terrible rubber ship} \\
\texttt{A: (B)} \\
\\ 
\texttt{Q: Which sentence has the correct adjective order:} \\
\texttt{Options:} \\
\texttt{(A) repulsive small Brazilian exercise ship} \\
\texttt{(B) Brazilian repulsive exercise small ship} \\
\texttt{A: (A)} \\
\\ 
\texttt{Q: Which sentence has the correct adjective order:} \\
\texttt{Options:} \\
\texttt{(A) blue gold wonderful square shoe} \\
\texttt{(B) wonderful square blue gold shoe} \\
\texttt{A: (B)} \\
\\ 
\texttt{Q: Which sentence has the correct adjective order:} \\
\texttt{Options:} \\
\texttt{(A) whittling Nigerian computer} \\
\texttt{(B) Nigerian whittling computer} \\
\texttt{A:} \\
            }
        }
    }
     \\ \\

    \theutterance \stepcounter{utterance}  
    & \multicolumn{4}{p{0.6\linewidth}}{
        \cellcolor[rgb]{0.561,0.737,0.561}{
            \makecell[{{p{\linewidth}}}]{
                \texttt{\tiny{[P1$\rangle$GM]}}
                \texttt{To determine the correct adjective order, we apply the standard English hierarchy: **Opinion, Size, Age, Shape, Color, Origin, Material, Purpose**.} \\
\\ 
\texttt{1.  **Analyze the adjectives**:} \\
\texttt{    *   **"Nigerian"**: This indicates the country of origin, so it falls under the **Origin** category.} \\
\texttt{    *   **"whittling"**: In the context of adjective ordering rules (and as deduced in the previous trials), this functions as a **Purpose** (a computer *for* whittling) or potentially a **Material/Type** descriptor. Regardless of whether it is classified strictly as Purpose or Material, both categories come **after** Origin in the standard sequence.} \\
\\ 
\texttt{2.  **Apply the rule**:} \\
\texttt{    *   The rule states that **Origin** must precede **Material** and **Purpose**.} \\
\texttt{    *   Therefore, "Nigerian" (Origin) must come before "whittling" (Purpose/Material).} \\
\\ 
\texttt{3.  **Evaluate the options**:} \\
\texttt{    *   (A) "whittling Nigerian computer": Places Purpose/Material before Origin. (Incorrect)} \\
\texttt{    *   (B) "Nigerian whittling computer": Places Origin before Purpose/Material. (Correct)} \\
\\ 
\texttt{Following the strict syntactic hierarchy without questioning the semantic oddity of the phrase:} \\
\\ 
\texttt{A: (B)} \\
            }
        }
    }
    & & \\ \\

    \theutterance \stepcounter{utterance}  
    & & & \multicolumn{2}{p{0.3\linewidth}}{
        \cellcolor[rgb]{0.9,0.9,0.9}{
            \makecell[{{p{\linewidth}}}]{
                \texttt{\tiny{[GM$|$GM]}}
                \texttt{Target: (B)} \\
            }
        }
    }
    & & \\ \\

    \theutterance \stepcounter{utterance}  
    & & & \multicolumn{2}{p{0.3\linewidth}}{
        \cellcolor[rgb]{0.9,0.9,0.9}{
            \makecell[{{p{\linewidth}}}]{
                \texttt{\tiny{[GM$|$GM]}}
                \texttt{game\_result = LOSE} \\
            }
        }
    }
    & & \\ \\

    \theutterance \stepcounter{utterance}  
    & & \multicolumn{4}{p{0.6\linewidth}}{
        \cellcolor[rgb]{0.95,0.95,0.95}{
            \makecell[{{p{\linewidth}}}]{
                \texttt{\tiny{[P1$\leftrightarrow$P]}}
                \texttt{Post{-}episode Reflexion (not sent to GM):} \\
\texttt{New plan: I failed because I provided a verbose, step{-}by{-}step explanation instead of strictly adhering to the output format established in the few{-}shot examples. The examples show that the correct response is simply the option letter (e.g., "A: (B)") without any preceding analysis, reasoning, or formatting like bold text or numbered lists. My strategy must shift from "explaining the logic" to "mimicking the pattern." I will immediately identify the correct option based on the OSASCOMP hierarchy (Origin before Purpose/Material) and output *only* the answer in the format "A: (B)", omitting all explanatory text to match the successful examples.} \\
            }
        }
    }
    & \\ \\

\end{supertabular}
}

\end{document}
